Esta sección presenta, a modo de referencia autocontenida, la formulación
matemática preliminar de las principales restricciones del modelo descritas
conceptualmente en el cuerpo del informe. Las expresiones corresponden a una
primera propuesta de modelación y podrán ser ajustadas durante las etapas
posteriores de implementación, verificación y validación.

\subsection*{D.1. Balance del inventario intermedio}

El inventario de botellas sin etiquetar debe conservar el flujo físico del
vino. Dado que el embotellado tiene un tiempo de disponibilidad de una etapa,
se plantea:

\begin{equation}
    s^{b}_{in}
    =
    s^{b}_{i,a(n)}
    +
    w^{b}_{i,a(n)}
    -
    \sum_{j\in J_i} w^{l}_{ijn},
    \qquad \forall i,n.
    \label{eq:anexoD-balance-intermedio}
\end{equation}

La producción sin etiquetar realizada en el nodo antecesor se encuentra
disponible en el nodo $n$, mientras que el etiquetado realizado en $n$ puede
consumir dicho inventario después de observar la demanda correspondiente al
nodo.

Adicionalmente, la cantidad etiquetada no puede superar el inventario de
botellas sin etiquetar efectivamente disponible:

\begin{equation}
    \sum_{j\in J_i} w^{l}_{ijn}
    \leq
    s^{b}_{i,a(n)}
    +
    w^{b}_{i,a(n)},
    \qquad \forall i,n.
    \label{eq:anexoD-disponibilidad-intermedia}
\end{equation}


\subsection*{D.2. Balance de producto terminado y pedidos atrasados}

Para cada combinación vino--etiqueta $(i,j)$, el inventario de producto
terminado y los pedidos atrasados deben satisfacer:

\begin{equation}
    s^{bl}_{ijn}
    -
    b^{bl}_{ijn}
    =
    s^{bl}_{ij,a(n)}
    -
    b^{bl}_{ij,a(n)}
    +
    w^{bl}_{ij,a(n)}
    +
    w^{l}_{ijn}
    -
    D_{ijn},
    \qquad \forall i,j,n.
    \label{eq:anexoD-balance-terminado}
\end{equation}

La producción acoplada disponible en el nodo $n$ corresponde a una decisión
realizada en su nodo antecesor, mientras que el etiquetado de inventario
intermedio puede reaccionar a la demanda observada en $n$. De esta forma, la
postergación permite desplazar parte de la diferenciación hacia un momento en
que se dispone de mayor información.


\subsection*{D.3. Capacidad de la línea}

La única línea disponible comparte su capacidad entre las operaciones de
embotellado, embotellado y etiquetado acoplados, etiquetado posterior y sus
respectivos \textit{set-ups}. Para cada nodo:

\begin{equation}
\begin{aligned}
    &0{,}00014\sum_i w^{b}_{in}
    +
    0{,}00028\sum_i\sum_{j\in J_i} w^{bl}_{ijn}
    +
    0{,}00028\sum_i\sum_{j\in J_i} w^{l}_{ijn}
    \\
    &\quad+
    1{,}5\sum_i z^{b}_{in}
    +
    1{,}5\sum_i\sum_{j\in J_i} z^{bl}_{ijn}
    +
    0{,}5\sum_i\sum_{j\in J_i} z^{l}_{ijn}
    \leq C,
    \qquad \forall n.
\end{aligned}
\label{eq:anexoD-capacidad}
\end{equation}

donde la capacidad disponible depende del escenario experimental considerado:

\begin{equation}
    C \in \{21,42,63,84\}.
    \label{eq:anexoD-niveles-capacidad}
\end{equation}

Esta restricción refleja que la postergación no constituye una fuente de
flexibilidad gratuita. Una botella que sigue la ruta desacoplada consume
capacidad durante el embotellado y vuelve a utilizarla posteriormente durante
el etiquetado. Por lo tanto, el beneficio de esperar nueva información debe
compensar el uso adicional de capacidad productiva.


\subsection*{D.4. Activación de \textit{set-ups}}

Las cantidades procesadas deben encontrarse asociadas a las variables binarias
que representan la activación de los \textit{set-ups} correspondientes. Se
plantean las siguientes relaciones:

\begin{equation}
    w^{b}_{in}
    \leq
    M^{b}(C)z^{b}_{in},
    \qquad \forall i,n,
    \label{eq:anexoD-setup-embotellado}
\end{equation}

\begin{equation}
    w^{bl}_{ijn}
    \leq
    M^{bl}(C)z^{bl}_{ijn},
    \qquad \forall i,j,n,
    \label{eq:anexoD-setup-acoplado}
\end{equation}

\begin{equation}
    w^{l}_{ijn}
    \leq
    M^{l}(C)z^{l}_{ijn},
    \qquad \forall i,j,n.
    \label{eq:anexoD-setup-etiquetado}
\end{equation}

Para evitar utilizar constantes $M$ excesivamente grandes, sus cotas se
derivan directamente de la capacidad disponible:

\begin{equation}
    M^{b}(C)
    =
    \max
    \left\{
        \frac{C-1{,}5}{0{,}00014},
        0
    \right\},
    \label{eq:anexoD-bigM-embotellado}
\end{equation}

\begin{equation}
    M^{bl}(C)
    =
    \max
    \left\{
        \frac{C-1{,}5}{0{,}00028},
        0
    \right\},
    \label{eq:anexoD-bigM-acoplado}
\end{equation}

\begin{equation}
    M^{l}(C)
    =
    \max
    \left\{
        \frac{C-0{,}5}{0{,}00028},
        0
    \right\}.
    \label{eq:anexoD-bigM-etiquetado}
\end{equation}

Estas cotas dependen del escenario de capacidad y evitan utilizar valores
arbitrariamente grandes de $M$ que puedan deteriorar innecesariamente la
relajación lineal del modelo.


\subsection*{D.5. Compatibilidad entre vinos y etiquetas}

El inventario de botellas sin etiquetar conserva flexibilidad únicamente entre
productos asociados al mismo vino base. Para la instancia analizada:

\begin{equation}
    J_1 = \{1,2,3\},
    \qquad
    J_2 = \{1,2\}.
    \label{eq:anexoD-compatibilidad}
\end{equation}

Por lo tanto, una botella correspondiente al Vino 1 puede transformarse en los
productos $(1,1)$, $(1,2)$ o $(1,3)$, mientras que una botella correspondiente
al Vino 2 puede asignarse únicamente a $(2,1)$ o $(2,2)$. En consecuencia, el
efecto de \textit{risk pooling} generado por la postergación ocurre entre
etiquetas que comparten el mismo vino base y no entre vinos diferentes.


\subsection*{D.6. Disponibilidad de estanques}

Los 20 estanques de 10.000 litros corresponden al vino preparado para alimentar
el proceso de embotellado y no representan capacidad de almacenamiento de
botellas sin etiquetar.

Sea $y_{in}$ el número entero de estanques preparados para el vino $i$ en el
nodo $n$, y $u_{in}$ su subutilización. El número de estanques utilizados debe
satisfacer:

\begin{equation}
    \sum_i y_{in}
    \leq
    20,
    \qquad \forall n.
    \label{eq:anexoD-numero-estanques}
\end{equation}

El volumen preparado se relaciona con la cantidad de botellas procesadas
mediante:

\begin{equation}
    10.000
    \left(
        y_{in}-u_{in}
    \right)
    =
    0{,}75
    \left(
        w^{b}_{in}
        +
        \sum_{j\in J_i} w^{bl}_{ijn}
    \right),
    \qquad \forall i,n.
    \label{eq:anexoD-balance-estanques}
\end{equation}

Dado que la subutilización máxima admisible de un estanque es del 100\%, no se
exige una utilización mínima de cada estanque. Sin embargo, el número total de
estanques continúa restringiendo físicamente el volumen de vino que puede
encontrarse preparado simultáneamente para alimentar la línea.


\subsection*{D.7. Condiciones iniciales y frontera temporal}

Los inventarios iniciales de botellas sin etiquetar, productos terminados y
pedidos atrasados son iguales a cero. La demanda correspondiente al nodo raíz
se satisface mediante actividades predefinidas de embotellado y etiquetado
equivalentes a la demanda de dicho nodo. Esta condición permite iniciar el
árbol desde un estado conocido y concentrar las decisiones de optimización en
la respuesta a la incertidumbre futura.

Asimismo, dado que el embotellado requiere una etapa para encontrarse
disponible, las decisiones adoptadas al final del horizonte no deben generar
artificialmente valor asociado a producto cuya disponibilidad se produzca
fuera de los tres períodos considerados.

El tratamiento de las variables terminales deberá revisarse en la formulación
definitiva, evitando incentivar producción cuya disponibilidad ocurra después
del tercer período, salvo que se incorpore explícitamente una condición de
valor residual.